% Einheit 4 von 4 – Gleichungen aufstellen (bank/lineare-gleichungen/e4.jsonl, zusammenbau v0.1)
%% TODO Zweigzeile Teil 1: was hier gelernt wird (Satz in Schülersprache aus dem Einheitstitel) – Einheit 4
\einheitenkopf[e4]{Einheit 4 von 4 · Gleichungen aufstellen}
\zweigzeile{ab Kl. 6, je nach Buch bis Kl. 8 (am Gymnasium ab Kl. 7) · P10}
\setcounter{aufgabe}{18}

%% TODO Titel: Ich-kann-Satz fehlt in der Bank (Platzhalter: Kettenname) – Nr. 19
%% TODO Anweisung: ein Satz über der ersten Teilaufgabe fehlt in der Bank – Nr. 19
\begin{aufgabe}{Aufstellen}
\begin{teile}
\teil „Eine Zahl plus 9 ergibt 23." – welche Gleichung passt? ($x$: die gesuchte Zahl) \\ \kreuz{$9x = 23$} \\ \kreuz{$x + 9 = 23$} \\ \kreuz{$x - 9 = 23$}
\end{teile}
\begin{teile}
\teil „Das Vierfache einer Zahl ist 36." – welche Gleichung passt? ($x$: die gesuchte Zahl) \\ \kreuz{$x + 4 = 36$} \\ \kreuz{$4x = 36$} \\ \kreuz{$x : 4 = 36$}
\end{teile}
\begin{teile}
\teil „Eine Zahl, vermindert um 7, ergibt 15." – welche Gleichung passt? ($x$: die gesuchte Zahl) \\ \kreuz{$7 - x = 15$} \\ \kreuz{$x - 7 = 15$} \\ \kreuz{$x + 7 = 15$}
\end{teile}
\begin{teile}
\teil „Die Hälfte einer Zahl ist 13." – welche Gleichung passt? ($x$: die gesuchte Zahl) \\ \kreuz{$2x = 13$} \\ \kreuz{$x : 2 = 13$} \\ \kreuz{$x - 2 = 13$}
\end{teile}
\begin{teile}
\teil „Das Doppelte einer Zahl, vermehrt um 5, ergibt 31." – welche Gleichung passt? ($x$: die gesuchte Zahl) \\ \kreuz{$2(x + 5) = 31$} \\ \kreuz{$5x + 2 = 31$} \\ \kreuz{$2x + 5 = 31$}
\end{teile}
\swfrage{Was bleibt gleich, was ändert sich?}
%% TODO Darstellung für schwach fehlt in der Bank (lineare-gleichungen-e4-k1-s2-v1; Abschnitt „Für schwache Schüler“)
\swz{Ich denke mir eine Zahl. Addiere ich 18, erhalte ich 45. Gleichung und Zahl? ($x$: die gesuchte Zahl)}{}
%% TODO Darstellung für schwach fehlt in der Bank (lineare-gleichungen-e4-k1-s3-v1; Abschnitt „Für schwache Schüler“)
\swz{Das Dreifache einer Zahl, vermehrt um 8, ergibt 44. Gleichung und Zahl? ($x$: die gesuchte Zahl)}{}
%% TODO Darstellung für schwach fehlt in der Bank (lineare-gleichungen-e4-k1-s4-v1; Abschnitt „Für schwache Schüler“)
\swz{Mia ist 4 Jahre älter als ihr Bruder. Zusammen sind sie 28 Jahre alt. Wie alt ist der Bruder? ($x$: Alter des Bruders in Jahren)}{}
%% TODO Darstellung für schwach fehlt in der Bank (lineare-gleichungen-e4-k1-s5-v1; Abschnitt „Für schwache Schüler“)
\swz[3]{Ein Rechteck hat den Umfang 46 cm, die Seite $a$ ist 8 cm lang. Wie lang ist die Seite $b$? Schreibe gegeben, gesucht, Formel, Rechnung.}{}
\swz[3]{Ein Freibad hatte am Samstag 380 Gäste, der Eintritt kostet 5 € pro Person. Der Betreiber sagt: „Mit $x$ Gästen mehr hätten wir 2\,400 € eingenommen." Stelle eine Gleichung auf und berechne, wie viele Gäste mehr es hätten sein müssen. \hfill (P10 2020 OS)}{}
\end{aufgabe}

%% TODO Titel: Ich-kann-Satz fehlt in der Bank (Platzhalter: Kettenname) – Nr. 20
%% TODO Anweisung: ein Satz über der ersten Teilaufgabe fehlt in der Bank – Nr. 20
\begin{aufgabe}{Formel umstellen}
%% TODO Darstellung für schwach fehlt in der Bank (lineare-gleichungen-e4-k2-s1-v1; Abschnitt „Für schwache Schüler“)
\swz{Flächeninhalt eines Rechtecks: $A = a \cdot b$ – nach $b$ umstellen.}{}
\end{aufgabe}

%% TODO Titel: Ich-kann-Satz fehlt in der Bank (Platzhalter: Kettenname) – Nr. 21
%% TODO Anweisung: ein Satz über der ersten Teilaufgabe fehlt in der Bank – Nr. 21
\begin{aufgabe}{Aufstellen – Fehler finden, Begründen, Darstellungswechsel, Anwendung}
\swz[3]{„Eine Zahl, um 6 vermindert, ergibt 10." Lisa schreibt $6 - x = 10$. Finde den Fehler.}{}
\swz[3]{„Das Vierfache der Summe aus einer Zahl und 3 ist 36." Tom schreibt $4x + 3 = 36$. Finde den Fehler.}{}
\swfrage{Für einen Ausflug wird berechnet, wie viele Busse nötig sind; die Gleichung liefert $x = 3{,}5$. Wie viele Busse werden bestellt? Begründe.}
\swz[3]{$x + 12 = 30$ – Zahlenrätsel dazu, und die Zahl?}{}
\swz[3]{Eine Taxifahrt kostet 4 € Grundpreis und 2,20 € je Kilometer. Frau Kaya zahlt 37 €. Wie lang war die Fahrt? ($x$: Strecke in km)}{}
\end{aufgabe}

\uebersichtskasten{Gleichung aus Text: Unbekannte festlegen (x = …), Text in Term übersetzen, Gleichung lösen, Antwort mit Einheit. \\ „Das Dreifache einer Zahl, vermindert um 4, ist 11": 3x − 4 = 11 \\ Formeln umstellen: wie eine Gleichung – nach der gesuchten Größe auflösen. \quad U = 2 · (a + b) → a = U : 2 − b}
